% !TEX root = ../dabirestan.tex
%% ============================================================
%%  فصل ۶ — قانون احتمال کل و قضیهٔ بیز
%% ============================================================
\chapter{قانون احتمال کل و قضیهٔ بیز}
\label{hs:ch6}

\chapterintro{\begin{itemize}
  \item با «شکستن مسئله به حالت‌ها» احتمال یک پیشامد را حساب کنید (قانون احتمال کل)؛
  \item با قضیهٔ بیز از «نتیجه» به «علت» برگردید؛
  \item بفهمید چرا یک آزمایش پزشکیِ «۹۵٪ دقیق» ممکن است بیشتر وقت‌ها اشتباه کند!
\end{itemize}}

%% ------------------------------------------------------------
\section{قانون احتمال کل}
%% ------------------------------------------------------------

گاهی احتمال یک پیشامد به این بستگی دارد که «کدام حالت» پیش آمده است. ایدهٔ قانون احتمال کل ساده است: \key{مسئله را به حالت‌های جدا از هم بشکن، هر حالت را جدا حساب کن و نتیجه‌ها را جمع بزن.}

\begin{example}[label=hs:ex:two-boxes]{دو جعبه}
جعبهٔ ۱ سه مهرهٔ قرمز و دو مهرهٔ آبی دارد و جعبهٔ ۲ یک مهرهٔ قرمز و چهار مهرهٔ آبی. سکه‌ای می‌اندازیم: اگر شیر آمد از جعبهٔ ۱ و اگر خط آمد از جعبهٔ ۲ یک مهره برمی‌داریم. احتمال اینکه مهره قرمز باشد چقدر است؟
\begin{solution}
نمودار درختی را رسم می‌کنیم:
\begin{figblock}
\begin{tikzpicture}[
  grow=left, level distance=3.2cm,
  level 1/.style={sibling distance=2.5cm},
  level 2/.style={sibling distance=1.15cm},
  every node/.style={font=\small},
  root/.style={circle, fill=primary, inner sep=2pt},
  st/.style={draw=primary!60, fill=primarylight, rounded corners=2pt, inner sep=3pt},
  red/.style={draw=accent, fill=notebg, rounded corners=2pt, inner sep=3pt},
  edge from parent/.style={draw=primary!55, thick},
  lab/.style={font=\scriptsize, fill=white, inner sep=1pt}]
  \node[root] {}
    child { node[st] {\rl{جعبهٔ ۱}}
      child { node[red] (r1) {\rl{قرمز}} edge from parent node[lab, above] {\lr{$\frac35$}} }
      child { node[st] {\rl{آبی}} edge from parent node[lab, below] {\lr{$\frac25$}} }
      edge from parent node[lab, above] {\lr{$\frac12$}} }
    child { node[st] {\rl{جعبهٔ ۲}}
      child { node[red] (r2) {\rl{قرمز}} edge from parent node[lab, above] {\lr{$\frac15$}} }
      child { node[st] {\rl{آبی}} edge from parent node[lab, below] {\lr{$\frac45$}} }
      edge from parent node[lab, below] {\lr{$\frac12$}} };
  \node[anchor=east, font=\small, text=accent!80!black] at ([xshift=-4mm]r1.west) {\lr{$\frac12\cdot\frac35=\frac3{10}$}};
  \node[anchor=east, font=\small, text=accent!80!black] at ([xshift=-4mm]r2.west) {\lr{$\frac12\cdot\frac15=\frac1{10}$}};
\end{tikzpicture}

\figcaption[hs:fig:two-boxes]{دو مسیر به «قرمز» می‌رسند؛ احتمال قرمز = جمع احتمال این دو مسیر.}
\end{figblock}
\[ \P(\fa{قرمز})=\frac12\cdot\frac35+\frac12\cdot\frac15=\frac3{10}+\frac1{10}=\frac25. \]
\end{solution}
\end{example}

\begin{definition}[label=hs:def:partition]{افراز}
پیشامدهای $B_1,B_2,\dots,B_k$ یک \term{افراز}{partition} از فضای نمونه‌اند اگر دوبه‌دو ناسازگار باشند و اجتماعشان کل $S$ باشد؛ یعنی در هر آزمایش \emph{دقیقاً یکی} از آن‌ها رخ دهد. (مثال بالا: «جعبهٔ ۱» و «جعبهٔ ۲».)
\end{definition}

\begin{theorem}[label=hs:thm:total]{قانون احتمال کل}
اگر $B_1,\dots,B_k$ افرازی از $S$ باشند، برای هر پیشامد $A$:
\[ \P(A)=\P(B_1)\,\P(A\mid B_1)+\P(B_2)\,\P(A\mid B_2)+\cdots+\P(B_k)\,\P(A\mid B_k). \]
\end{theorem}

\textbf{چرا؟} پیشامد $A$ به تکه‌های جدا از هم $A\cap B_1,\ A\cap B_2,\dots$ شکسته می‌شود. طبق قانون جمع، احتمال $A$ جمع احتمال این تکه‌هاست و طبق قانون ضرب $\P(A\cap B_i)=\P(B_i)\P(A\mid B_i)$. در نمودار درختی، این یعنی «جمع احتمال همهٔ مسیرهایی که به $A$ می‌رسند».

\begin{example}[label=hs:ex:factory]{کارخانه}
در یک کارخانه، دستگاه $M_1$ شصت درصد و دستگاه $M_2$ چهل درصد لامپ‌ها را تولید می‌کند. ۲٪ لامپ‌های $M_1$ و ۵٪ لامپ‌های $M_2$ معیوب‌اند. لامپی به تصادف انتخاب می‌شود. احتمال معیوب بودن آن چقدر است؟
\begin{solution}
$\P(\fa{معیوب})=0.6\times0.02+0.4\times0.05=0.012+0.020=0.032$؛ یعنی ۳٫۲٪.
\end{solution}
\end{example}

\begin{example}[label=hs:ex:bus]{دیر رسیدن سرویس}
احتمال بارانی بودن یک روز $0.3$ است. در روزهای بارانی سرویس مدرسه با احتمال $0.5$ و در روزهای غیربارانی با احتمال $0.1$ دیر می‌رسد. احتمال دیر رسیدن سرویس در یک روز دلخواه چقدر است؟
\begin{solution}
$\P(\fa{دیر})=0.3\times0.5+0.7\times0.1=0.15+0.07=0.22$.
\end{solution}
\end{example}

%% ------------------------------------------------------------
\section{قضیهٔ بیز: برگشت از نتیجه به علت}
%% ------------------------------------------------------------

در مثال دو جعبه پرسیدیم: «احتمال قرمز بودن مهره چقدر است؟» حالا سؤال را برعکس می‌کنیم: \key{مهره قرمز درآمده؛ احتمال اینکه از جعبهٔ ۱ برداشته شده باشد چقدر است؟} این یعنی برگشت از \emph{نتیجه} (قرمز) به \emph{علت} (کدام جعبه).

\begin{theorem}[label=hs:thm:bayes]{قضیهٔ بیز}
اگر $B_1,\dots,B_k$ افرازی از $S$ باشند و $\P(A)>0$، آنگاه برای هر $i$:
\[ \P(B_i\mid A)=\frac{\P(B_i)\,\P(A\mid B_i)}{\P(A)}, \]
که در آن مخرج، $\P(A)$، با قانون احتمال کل حساب می‌شود.
\end{theorem}

\textbf{چرا؟} طبق قانون ضرب، $\P(A\cap B_i)$ را به دو روش می‌نویسیم: $\P(B_i)\P(A\mid B_i)=\P(A)\P(B_i\mid A)$. با تقسیم بر $\P(A)$ قضیه به دست می‌آید.

\begin{intuition}
در نمودار درختی: احتمال اینکه از مسیر $B_i$ آمده باشیم، به شرط اینکه به $A$ رسیده‌ایم، برابر است با
\[ \frac{\text{\rl{احتمال مسیری که از $B_i$ به $A$ می‌رسد}}}{\text{\rl{جمع احتمال همهٔ مسیرهایی که به $A$ می‌رسند}}}. \]
\end{intuition}

\begin{example}[label=hs:ex:boxes-bayes]{دو جعبه، برعکس}
در مثال \ref{hs:ex:two-boxes} مهرهٔ برداشته‌شده قرمز است. احتمال اینکه از جعبهٔ ۱ باشد چقدر است؟
\begin{solution}
از نمودار درختی: مسیر «جعبهٔ ۱، قرمز» احتمال $\tfrac3{10}$ و کل مسیرهای قرمز $\tfrac25$ است. پس
\[ \P(\fa{جعبهٔ ۱}\mid\fa{قرمز})=\frac{3/10}{2/5}=\frac34. \]
قبل از دیدن مهره، احتمال جعبهٔ ۱ برابر $\tfrac12$ بود؛ قرمز بودن مهره (که در جعبهٔ ۱ رایج‌تر است) آن را به $\tfrac34$ رساند.
\end{solution}
\end{example}

\begin{example}[label=hs:ex:factory-bayes]{کدام دستگاه؟}
در مثال \ref{hs:ex:factory}، لامپ انتخاب‌شده معیوب است. احتمال اینکه دستگاه $M_2$ آن را ساخته باشد چقدر است؟
\begin{solution}
$\P(M_2\mid\fa{معیوب})=\dfrac{0.4\times0.05}{0.032}=\dfrac{0.020}{0.032}=0.625$. با اینکه $M_2$ فقط ۴۰٪ لامپ‌ها را می‌سازد، ۶۲٫۵٪ لامپ‌های معیوب مال آن است؛ چون درصد خرابی‌اش بیشتر است.
\end{solution}
\end{example}

\begin{example}[label=hs:ex:bus-bayes]{باران و سرویس}
در مثال \ref{hs:ex:bus}، سرویس امروز دیر آمده است. احتمال اینکه امروز بارانی باشد چقدر است؟
\begin{solution}
$\P(\fa{بارانی}\mid\fa{دیر})=\dfrac{0.3\times0.5}{0.22}=\dfrac{0.15}{0.22}\approx0.68$. دیر آمدن سرویس، احتمال باران را از ۳۰٪ به حدود ۶۸٪ رساند.
\end{solution}
\end{example}

\subsection{آزمایش پزشکی: یک نتیجهٔ شگفت‌انگیز}

\begin{example}[label=hs:ex:medical]{آزمایش ۹۵٪ دقیق}
یک بیماری در ۱٪ مردم وجود دارد. آزمایشی برای آن داریم که:
اگر فرد بیمار باشد، با احتمال ۹۵٪ جواب مثبت می‌دهد؛ و اگر فرد سالم باشد، با احتمال ۵٪ به اشتباه جواب مثبت می‌دهد.
آزمایش کسی مثبت شده است. احتمال اینکه واقعاً بیمار باشد چقدر است؟
\begin{solution}
بیشتر مردم حدس می‌زنند حدود ۹۵٪. بیایید با \emph{تعداد آدم‌ها} حساب کنیم. فرض کنید ۱۰۰۰۰ نفر آزمایش شوند:
\begin{figblock}
\begin{tikzpicture}[
  box/.style={draw=#1, fill=#1!10, rounded corners=3pt, minimum width=2.3cm, minimum height=0.8cm, font=\small, align=center},
  arr/.style={-{Stealth[length=2mm]}, thick, graytext}]
  \node[box=primary] (all) at (0,0) {\rl{۱۰۰۰۰ نفر}};
  \node[box=warncolor] (d) at (3.4,-1.5) {\rl{۱۰۰ بیمار}};
  \node[box=excolor] (h) at (-3.4,-1.5) {\rl{۹۹۰۰ سالم}};
  \node[box=warncolor] (dp) at (5,-3) {\rl{۹۵ مثبت}};
  \node[box=graytext] (dn) at (1.8,-3) {\rl{۵ منفی}};
  \node[box=accent] (hp) at (-1.8,-3) {\rl{۴۹۵ مثبت}};
  \node[box=graytext] (hn) at (-5,-3) {\rl{۹۴۰۵ منفی}};
  \draw[arr] (all) -- node[above right, font=\scriptsize]{\lr{$1\%$}} (d);
  \draw[arr] (all) -- node[above left, font=\scriptsize]{\lr{$99\%$}} (h);
  \draw[arr] (d) -- node[right, font=\scriptsize]{\lr{$95\%$}} (dp);
  \draw[arr] (d) -- (dn);
  \draw[arr] (h) -- node[right, font=\scriptsize]{\lr{$5\%$}} (hp);
  \draw[arr] (h) -- (hn);
\end{tikzpicture}

\figcaption[hs:fig:medical]{از ۱۰۰۰۰ نفر: ۹۵ بیمارِ مثبت و ۴۹۵ سالمِ مثبت.}
\end{figblock}
روی هم $95+495=590$ نفر جواب مثبت گرفته‌اند که فقط ۹۵ نفرشان واقعاً بیمارند:
\[ \P(\fa{بیمار}\mid\fa{مثبت})=\frac{95}{590}\approx0.16. \]
\textbf{فقط حدود ۱۶٪!} دلیلش این است که بیماری \emph{نادر} است: آن‌قدر آدم سالم زیاد است که حتی ۵٪ اشتباه در میان آن‌ها، از تعداد کل بیماران بیشتر می‌شود. با فرمول بیز هم همین به دست می‌آید:
\[ \P(\fa{بیمار}\mid\fa{مثبت})=\frac{0.01\times0.95}{0.01\times0.95+0.99\times0.05}=\frac{0.0095}{0.059}\approx0.16. \]
\end{solution}
\end{example}

\begin{warning}
$\P(\fa{مثبت}\mid\fa{بیمار})=0.95$ با $\P(\fa{بیمار}\mid\fa{مثبت})\approx0.16$ خیلی فرق دارد. اشتباه گرفتن این دو یکی از رایج‌ترین خطاها، حتی بین پزشکان، است. به همین دلیل پزشکان معمولاً آزمایش مثبت را با آزمایش دوم تأیید می‌کنند.
\end{warning}

\begin{moreinfo}{مسئلهٔ مونتی هال}
در یک مسابقه پشت یکی از سه در یک خودرو و پشت دو در دیگر بز است. شما در ۱ را انتخاب می‌کنید. مجری که می‌داند خودرو کجاست، یکی از دو در دیگر را که پشتش بز است باز می‌کند (مثلاً در ۳) و می‌پرسد: «می‌خواهید انتخابتان را به در ۲ عوض کنید؟»

جواب شگفت‌انگیز: \textbf{بله! با عوض کردن، احتمال برد از $\tfrac13$ به $\tfrac23$ می‌رسد.} دلیل ساده: انتخاب اول شما با احتمال $\tfrac23$ اشتباه است (خودرو پشت یکی از دو در دیگر است). در این حالت، مجری مجبور است تنها درِ بزدارِ باقی‌مانده را باز کند؛ پس در دیگری که باز نشده، حتماً خودرو دارد. یعنی با احتمال $\tfrac23$، عوض کردن شما را برنده می‌کند. (این را با قضیهٔ بیز هم می‌توان دقیق ثابت کرد.)
\end{moreinfo}

%% ------------------------------------------------------------
\section*{خلاصهٔ فصل ۶}
\addcontentsline{toc}{section}{خلاصه و تمرین‌ها}
%% ------------------------------------------------------------

\begin{summary}[title={\leavevmode\boxicon{accent}{$\Sigma$}\ \ خلاصهٔ فصل}]
\begin{itemize}
  \item \textbf{افراز:} پیشامدهای ناسازگاری که دقیقاً یکی‌شان رخ می‌دهد.
  \item \textbf{قانون احتمال کل:} $\P(A)=\sum_i\P(B_i)\P(A\mid B_i)$؛ در نمودار درختی: جمع مسیرهایی که به $A$ می‌رسند.
  \item \textbf{قضیهٔ بیز:} $\P(B_i\mid A)=\dfrac{\P(B_i)\P(A\mid B_i)}{\P(A)}$؛ برگشت از نتیجه به علت.
  \item برای مسائل بیز، فکر کردن با «تعداد آدم‌ها» (مثلاً ۱۰۰۰۰ نفر) خیلی کمک می‌کند.
  \item $\P(A\mid B)\ne\P(B\mid A)$؛ نادر بودن یک حالت (مثل یک بیماری) نتیجه را خیلی تغییر می‌دهد.
\end{itemize}
\end{summary}

\subsection*{تمرین‌ها}

\begin{exercise}
جعبهٔ ۱ دو مهرهٔ سفید و سه مهرهٔ سیاه و جعبهٔ ۲ چهار مهرهٔ سفید و یک مهرهٔ سیاه دارد. تاسی می‌اندازیم: اگر ۱ یا ۲ آمد از جعبهٔ ۱ و در غیر این صورت از جعبهٔ ۲ یک مهره برمی‌داریم. (الف) احتمال سفید بودن مهره چقدر است؟ (ب) اگر مهره سفید باشد، احتمال اینکه از جعبهٔ ۱ باشد چقدر است؟
\answer{(الف) $\tfrac13\cdot\tfrac25+\tfrac23\cdot\tfrac45=\tfrac{10}{15}=\tfrac23$ \quad (ب) $\tfrac{2/15}{2/3}=\tfrac15$}
\end{exercise}

\begin{exercise}
۶۰٪ دانش‌آموزان یک مدرسه دختر و ۴۰٪ پسرند. ۳۰٪ دخترها و ۵۰٪ پسرها عضو تیم ورزشی‌اند. (الف) چند درصد دانش‌آموزان عضو تیم‌اند؟ (ب) یک عضو تیم به تصادف انتخاب می‌شود؛ احتمال اینکه دختر باشد چقدر است؟
\answer{(الف) ۳۸٪ \quad (ب) $\tfrac{0.18}{0.38}=\tfrac9{19}\approx0.47$}
\end{exercise}

\begin{exercise}
سه دستگاه $A$، $B$ و $C$ به ترتیب ۵۰٪، ۳۰٪ و ۲۰٪ قطعات یک کارخانه را می‌سازند و درصد خرابی آن‌ها به ترتیب ۱٪، ۲٪ و ۳٪ است. (الف) احتمال معیوب بودن یک قطعه چقدر است؟ (ب) اگر قطعه‌ای معیوب باشد، احتمال اینکه دستگاه $C$ آن را ساخته باشد چقدر است؟
\answer{(الف) $0.017$ \quad (ب) $\tfrac{0.006}{0.017}=\tfrac6{17}\approx0.35$}
\end{exercise}

\begin{exercise}
یک بیماری در ۲٪ مردم وجود دارد. آزمایشی برای بیماران با احتمال ۹۰٪ و برای افراد سالم با احتمال ۱۰٪ جواب مثبت می‌دهد. با روش «۱۰۰۰ نفر» احتمال بیمار بودن کسی را که آزمایشش مثبت شده حساب کنید.
\answer{۱۸ بیمارِ مثبت و ۹۸ سالمِ مثبت: $\tfrac{18}{116}\approx0.16$}
\end{exercise}

\begin{exercise}
یک فیلتر ایمیل ۹۰٪ هرزنامه‌ها را تشخیص می‌دهد و ۲٪ ایمیل‌های معمولی را هم به اشتباه هرزنامه تشخیص می‌دهد. اگر ۲۰٪ ایمیل‌ها هرزنامه باشند، احتمال اینکه ایمیلی که فیلتر آن را هرزنامه تشخیص داده واقعاً هرزنامه باشد چقدر است؟
\answer{$\tfrac{0.2\times0.9}{0.2\times0.9+0.8\times0.02}=\tfrac{0.18}{0.196}\approx0.92$}
\end{exercise}
